%%%%%%%%%%%%%%%%%%%%%%%%%%%%%%%%%%%%%%%%%%%%%%%%%%%%%%%%%%%%%%%%%%%%%%
% Lớp 10 | Thống kê | Số gần đúng. Sai số
% Dạng 0D6?1-4: Viết số gần đúng cho số đúng biết độ chính xác
% Số câu: 14  (sinh tự động bởi tach_cau.py; từ nay thầy sửa trực tiếp file này)
%%%%%%%%%%%%%%%%%%%%%%%%%%%%%%%%%%%%%%%%%%%%%%%%%%%%%%%%%%%%%%%%%%%%%%

%%%%%%%%%%%%%%%% TRẮC NGHIỆM 4 PHƯƠNG ÁN %%%%%%%%%%%%%%%%

% nguồn: MaHoaNH_TracNghiem_Lop10
\begin{ex}%[0D6N1-4]%[ma:99d4f190]
	Cho giá trị gần đúng của $\pi $ là $ a=3,141592653589$ với  độ chính xác $10^{-10}$. Hãy viết số quy tròn của số $a$.
	 \choice
		{$  a=3{,}1415926535$}
		{$ a=3{,}1415926536$} 
		{$ a=3{,}141592653$} 
		{\True $ a=3{,}141592654$} 
	\loigiai{ Với độ chính xác $10^{-10}$ thì $a$ quy tròn đến vị trí số thứ  $10^{-9}$. Khi đó \\
	$ a=3{,}141592653589 \approx  a=3{,}141592654$.}
\end{ex}

% nguồn: MaHoaNH_TracNghiem_Lop10
\begin{ex}%[0D6N1-4]%[ma:47518243]
	Số quy tròn của  số gần đúng $ a=35{,}675$ với độ chính xác $ 0{,}02$ là
	 \choice
		 {$35{,}6$}
		{$35{,}67$} 
		 {$36$}
		{\True  $35{,}7$}
	\loigiai{ Với độ chính xác $ 0{,}02$ thì $a$  quy tròn đến phần chục. Khi đó $a=35{,}675 \approx 35{,}7$.
	 }
\end{ex}

% nguồn: MaHoaNH_TracNghiem_Lop10
\begin{ex}%[0D6N1-4]%[ma:bd32f84f]
	Trong một cuộc điều tra dân số, người ta báo dân số của tỉnh A là $ 1\, 279 \,  425\pm 300$ người. Hãy viết số quy tròn số dân trên.
	 \choice
		{$ 1\, 270\, 000$ người} 
		{$ 1\, 279\, 400$ người} 
		{$  1\, 280\, 000$ người}
		{\True $  1\, 279\, 000$ người}
	\loigiai{ $d=300$ nên  dân số tỉnh A quy tròn đến hàng nghìn. Khi đó
	$1\, 279\,  425\pm 300 \approx 1\, 279\, 000$.}
\end{ex}

% nguồn: MaHoaNH_TracNghiem_Lop10
\begin{ex}%[0D6N1-4]%[ma:e08b4cd2]
	Đo độ  cao một ngọn cây là $h=17{,}14\,\text{m}\pm 0{,}3\,\text{m}$. Hãy viết số quy tròn của số $17{,}14$?
	 \choice
		{$17{,}1$} 
		{$17{,}15$} 
		 {$17{,}2$}
		{\True  $17$}
	 \loigiai{ Do độ chính xác là $0{,}3$ nên $h$ quy tròn đến hàng đơn vị, do đó $h=17{,}14 \approx 17$.
	 }
\end{ex}

% nguồn: MaHoaNH_TracNghiem_Lop10
\begin{ex}%[0D6N1-4]%[ma:bf6153fe]
	Cho  số $\overline{a}=4{,}1356\pm 0{,}001$. Số quy tròn của số gần đúng $ 4{,}1356$ là
	\choice 
		{$  4{,}135$}
		{$  4{,}13$}
		 {$ 4{,}136$}
		{\True $ 4{,}14$} 
	 \loigiai{
	Vì độ chính xác đến hàng phần nghìn (độ chính xác là $ 0{,}001$) nên ta quy tròn số $ 4{,}1356$  đến hàng phần phần trăm theo quy tắc làm tròn. Vậy số quy tròn của số $ 4{,}1356$ là $ 4{,}14$.}
\end{ex}

% nguồn: MaHoaNH_TracNghiem_Lop10
\begin{ex}%[0D6N1-4]%[ma:fae125a1]
	Cho số gần đúng $a=8\,141\,378$ với độ chính xác $d=300$. Hãy viết quy tròn số  $a$.
	 \choice
		{$8\,141\,400$} 
		{$8\,142\,400$} 
		 {\True $8\,141\,000$}
		 {$8\,141\,300$}
	\loigiai{ $d=300$ nên $a$ quy tròn đến hàng  nghìn. Khi đó $a=8\,141\,378 \approx 8\,141\,000$.
	 }
\end{ex}

% nguồn: MaHoaNH_TracNghiem_Lop10
\begin{ex}%[0D6N1-4]%[ma:b446fdfd]
	Giá trị gần đúng của $\sqrt{3}$ chính xác đến  hàng phần trăm là
	 \choice
		{\True $1{,}73$} 
		{$2$} 
		{$1{,}7$} 
		{$1{,}732$} 
	\loigiai{ 
		Giá trị gần đúng của  $\sqrt{3}$ chính xác đến hàng phần trăm là $1{,}73$.
	 }
\end{ex}

% nguồn: MaHoaNH_TracNghiem_Lop10
\begin{ex}%[0D6H1-4]%[ma:30d059cf]
	Theo thống kê, dân số Việt Nam năm $ 2002$ là $ 79\, 715 \, 675$ người. Giả sử sai số tuyệt đối của số liệu thống kê này nhỏ hơn $ 10\, 000$ người. Hãy viết  số quy tròn của số trên.
	 \choice
		{$ 79\, 710\, 000$ người} 
		{$  79\, 716\, 000$ người}
		{\True $ 79\, 720\, 000$  người}
		{$ 79\, 000\, 000$ người} 
	\loigiai{ 
	Do sai số tuyệt đối của thống kê này nhỏ hơn $ 10\, 000$ người nên ta  quy tròn số trên đến hàng chục nghìn. Vậy số quy tròn của số $ 79 \, 715 \, 675$ là $ 79\, 720\, 000$.}
\end{ex}


%%%%%%%%%%%%%%%% ĐÚNG/SAI %%%%%%%%%%%%%%%%

% nguồn: MaHoaNH_DungSai_Lop10
\begin{ex}%[0D6H1-4]%[ma:70672e47]
	Một  bao gạo có ghi $12 \pm 0{,}2$ kg.
	Các khẳng định  sau đây đúng hay sai?
	\choiceTF 
		{\True Khối lượng của bao gạo không nhỏ hơn $11,8$  kg}
		{\True Sai số tuyệt đối không vượt  qua $0{,}2$ kg}
		{\True Số  quy tròn khối lượng của bao gạo là $12$ kg}
		{Sai  số tương đối không vượt quá $1{,}5\%$}
	 \loigiai{
		\begin{itemchoice}
			\itemch Khối  lượng của bao gạo không nhỏ hơn $12-0,2 = 11,8$ kg
			 \itemch Sai số tuyệt đối không vượt qua $0{,}2$ kg
			\itemch Số quy tròn khối lượng của bao gạo là $12$  kg
			\itemch Sai số tương đối không vượt quá $\dfrac{0,2}{12} \approx  1{,}67\%$.
		\end{itemchoice}
	} 
\end{ex}

% nguồn: MaHoaNH_DungSai_Lop10
\begin{ex}%[0D6H1-4]%[ma:4eb1d4c8]
	Cho  số gần đúng $a=2{,}82842$ với độ chính xác $0{,}0005$. \\
	Các  khẳng định sau đây đúng hay sai?
	 \choiceTF
		{\True Sai số tuyệt đối không vượt quá $0{,}0005$} 
		{\True Sai số tương đối không vượt quá $0{,}016\%$} 
		{Giá trị đúng $\overline{a}$ không  nhỏ hơn $2{,}8285$}
		{\True Số quy tròn của $a$ là  $2{,}828$}
	 \loigiai{
		\begin{itemchoice}
			\itemch Sai số tuyệt đối không vượt quá  $0{,}0005$
			\itemch Sai số tương đối không vượt  quá $\dfrac{0,005}{2,82841} \approx 0{,}0176 \%$.
			\itemch Giá trị đúng $\overline{a}$ không nhỏ hơn $2{,}82842- 0{,}0005=2{,}82792$. 
			\itemch Số  quy tròn của $a$ là $2{,}828$.
		\end{itemchoice}
	} 
\end{ex}

% nguồn: MaHoaNH_DungSai_Lop10
\begin{ex}%[0D6V1-4]%[ma:96bd7008]
	Một công ty sử dụng dây chuyền A để đóng gạo  vào bao với khối lượng mong muốn $5$ kg. Trên bao bì ghi thông tin khối lượng $5\pm 0,2$ kg. Gọi $\bar{a}$ là khối lượng thực tế của một bao gạo do dây chuyền $A$ đóng gói. Khi đó
	\choiceTF 
		{\True  Số đúng $\bar{a}$}
		{$5{,}2$  là số gần đúng của $\bar{a}$}
		{Độ chính xác là $ d=5$  kg}
		{\True Giá trị  của $\bar{a}$ nằm trong đoạn $[4,8;5,2]$}
	 \loigiai{
		\begin{itemchoice}
			\itemch Đúng. Vì Số đúng $\bar{a}$, cũng  là khối lượng thực tế của bao gạo (chưa biết chính xác $\bar{a}$).
			\itemch Sai. Vì Khối lượng mong muốn là $5$ kg nên số $5$ là số gần đúng của  $\bar{a}$.
			\itemch Sai. Vì Độ chính xác là $d=0,2$ kg. 
			\itemch Đúng. Vì Giá  trị $\bar{a}$ nằm trong đoạn $[5-0,2; 5+0,2]=[4,8; 5,2]$.
		\end{itemchoice}
	 }
\end{ex}

% nguồn: MaHoaNH_DungSai_Lop10
\begin{ex}%[0D6V1-4]%[ma:c7354cf3]
	Cho ba giá trị gần  đúng của $\dfrac{3}{7}$ là $0{,}429$; $0{,}4$ và $0{,}42$. Xét tính đúng, sai của các mệnh đề sau
	 \choiceTF
		{\True Công thức đánh giá sai số tuyệt  đối là $\Delta=|\overline{a} - a|$}
		{\True Xét số gần đúng $0{,}429$ ta  có $\Delta_1=\left|\dfrac{3}{7} - 0{,}429\right|<0{,}0005$}
		{\True Xét số gần đúng  $0{,}4$ ta có $\Delta_2=\left|\dfrac{3}{7} - 0{,}4\right|<0{,}03$}
		{\True Xét số  gần đúng $0{,}42$ ta có $\Delta_2=\left|\dfrac{3}{7} - 0{,}42\right|<0{,}009$}
	 \loigiai{
		\begin{itemchoice}
			 \itemch Đúng. Ta sử dụng công thức đánh giá sai số tuyệt đối là $\Delta=|\overline{a} - a|$.
			\itemch Đúng. Xét số gần đúng $0{,}429$ ta có $\Delta_1=\left|\dfrac{3}{7} -  0{,}429\right|<0{,}0005$.
			\itemch  Đúng. Xét số gần đúng $0{,}4$ ta có $\Delta_2=\left|\dfrac{3}{7} - 0{,}4\right|<0{,}03$.
			\itemch Đúng. Xét số gần đúng $0{,}42$ ta có $\Delta_2=\left|\dfrac{3}{7}  - 0{,}42\right|<0{,}009$.
		\end{itemchoice}
	} 
\end{ex}

% nguồn: MaHoaNH_DungSai_Lop10
\begin{ex}%[0D6V1-4]%[ma:1ca7e6ec]
	Xét tính đúng, sai của các mệnh đề sau 
	\choiceTF 
		{Cho $\overline{a}=\sqrt{7}=2{,}645751$. Số gần đúng của $\overline{a}$ với độ  chính xác $d=0{,}004$ là $2{,}6$}
		{Cho $\overline{a}=\dfrac{\sqrt[3]{2}}{4}=0{,}314980\ldots$ Số gần đúng  của $\overline{a}$ với độ chính xác $d=0{,}01$ là $0{,}315$}
		{\True Cho số gần đúng  $a=1\,624\,192$ với độ chính xác $d=300$. Số quy tròn của $a$ là $1\,624\,000$}
		{\True Viết quy tròn số gần đúng $a=3{,}26356$ biết $\overline{a}=3{,}26356\pm 0{,}001$  là $3{,}26$}
	\loigiai{ 
		\begin{itemchoice}
			\itemch Sai.  Kết quả là $2{,}65$.
			\itemch Sai.  Kết quả là $0{,}3$.
			\itemch Đúng. Số  quy tròn của $a$ là $1\,624\,000$.
			\itemch Đúng. $\overline{a}=3{,}26356\pm 0{,}001\Rightarrow$ Độ chính xác đến hàng phần nghìn (độ chính  xác là $d=0{,}001$).\\
			Ta quy  tròn đến hàng phần trăm. Vậy số quy tròn của $a$ là $3{,}26$.
		\end{itemchoice}
	} 
\end{ex}


%%%%%%%%%%%%%%%% TRẢ LỜI NGẮN %%%%%%%%%%%%%%%%

% nguồn: MaHoaNH_TraLoiNgan_Lop10
\begin{bt}%[0D6H1-4]%[ma:04ead347]
	Trên bao bì của một sản phẩm có ghi \lq\lq khối lượng tịnh $200 \pm 2$ gam\,\rq\rq. Biết khối lượng đúng của bao bì sản phẩm đó  thuộc đoạn $[m; n]$, với $m$; $n$ là các số tự nhiên. Tính $S=m+n$.
	\shortans{400} 
	 \loigiai{
		Khối  lượng đúng của bao sản phẩm $\bar{a}$ (tính theo gam) thuộc đoạn $[198; 202]$.\\
	Khi đó $\heva{&m=198\\&  n=202} \Rightarrow S=400.$}
\end{bt}
